\section{Training}
\label{sec:training}

\subsection{Stages}
We adopt a three-stage training recipe:

\begin{enumerate}
  \item \textbf{Stage 1 -- Alignment.} The audio encoder and LLM are frozen;
        only the adapter is trained on (audio, transcript) pairs to align
        speech representations with the LLM token space.
  \item \textbf{Stage 2 -- Supervised fine-tuning (SFT).} The adapter and the
        LLM are jointly fine-tuned on a balanced multi-task mixture
        (transcription, contextual ASR, instruction-style ASR).
  \item \textbf{Stage 3 -- Preference optimization.} We further refine the
        model with DPO/RLAIF using preference pairs sampled from Stage 2.
\end{enumerate}

\subsection{Objective}
% TODO: Token-level cross-entropy on the transcript; loss masking on prompt tokens.

\subsection{Optimization}
% TODO: AdamW, lr schedule (warmup + cosine), batch size in tokens, sequence length.

\begin{table}[t]
  \caption{Training hyperparameters per stage.}
  \label{tab:hparams}
  \centering
  \small
  \begin{tabular}{lccc}
    \toprule
    & Stage 1 & Stage 2 & Stage 3 \\
    \midrule
    Trainable params & adapter & adapter+LLM & adapter+LLM \\
    Learning rate    & 1e-3    & 2e-5        & 5e-7 \\
    Batch (audio hr) & ?       & ?           & ? \\
    Steps            & ?       & ?           & ? \\
    Precision        & bf16    & bf16        & bf16 \\
    \bottomrule
  \end{tabular}
\end{table}

\subsection{Infrastructure}
% TODO: GPUs/TPUs, parallelism strategy (FSDP / DeepSpeed-Zero), throughput.
